\justifying 
The Chapter presents a detailed examination of the Virtual Lab Assistant project.
Analysis simplifies the project by dividing it into different parts, which is a crucial process for effective software development.
\item The organization of this Chapter is as follows:
Section 3.1 focuses on Requirement Collection and Identification.
The Software Requirement Specification (SRS) is presented in Section 3.2.
Section 3.3 summarizes the entire Chapter.


\section{Requirement Collection and Identification}
Requirement Collection and Identification is a critical step in the development of the Virtual Lab Assistant system. It involves gathering, analyzing, and defining the needs of users, including students and administrators.

The requirements were identified through observation of existing practical examination systems and understanding the challenges faced during manual practical sessions. Key requirements include user authentication, practical execution, real-time query evaluation, performance tracking, and admin control functionalities.


\section{Software Requirement Specification (SRS)}
The Software Requirement Specification (SRS) for the Virtual Lab Assistant defines its functionality, user roles, and environmental needs. It ensures a clear understanding between developers, faculty, and students, minimizing ambiguity during development.
This section outlines the functional, non-functional, and external requirements essential to the project.
\subsection{Product Features}
- User Management: Registration and login system for students using PRN numbers, and secure login for admin users.  

\item- Subject & Practical Management:
: Supports multiple subjects (DBMS, C, C++, Python). Each subject contains 10 practicals with 10 related questions 
\item - Exam Controls: admin can see the progress of students and modify questions per subject  Students can only attempt practicals on their scheduled date.

\item- Question Handling:: Each question includes an example, a hint, and code/SQL input support  

\item- Evaluate & Submit Feature: Students can simulate “evaluate” operations and submit their solutions for evaluation

\item- Timer & Attempt Control: Each exam runs under a timer (30 minutes) with automatic submission on timeout. Login attempts are limited to prevent misuse
\item- Feedback & Help Section: Students can send feedback and contact administrators directly from the platform.
\item- Performance Tracking: Admin can monitor completion percentages of subjects and practicals through progress indicators.

\subsection{Operating Environment}
Software and hardware environments, referred to as the operating environment for the project, are outlined as follows:


\begin{enumerate}
  \item	Operating Environment:
The Virtual Lab Assistant is a web-based platform accessible from any modern browser such as Google Chrome, Firefox, or Edge.
It can be hosted on any web server supporting HTML, CSS, and JavaScript for the frontend and future integration with PHP or Node.js for the backend.
Secure communication (HTTPS) ensures safe data transfer.


 
  \item	Hardware Environment:
The system requires only basic hardware — desktop, laptop, or tablet with an internet connection.
A server with at least 4 GB RAM and SSD storage is recommended for efficient performance and scalability.




\end{enumerate}
\newpage
\subsection{Assumptions} 
The following assumptions are made during the development and usage of the Virtual Lab Assistant:

\begin{itemize}

    \item	Users (students and administrators) will have access to stable internet connectivity.
    \item	Students have basic computer and internet literacy.
    \item	Admin users will maintain the correctness of practical questions 
    \item	All users will operate the system through authorized login credentials.
    \item	The server environment will support standard web technologies and ensure data backup.
    \item	The system will be scalable to accommodate an increasing number of students and subjects.

 
\end{itemize}

\subsection{ Functional Requirements}
The Virtual Lab Assistant system includes the following functional requirements:
\subsubsection{\textbf{System}}
\begin{enumerate}
  

   \item 	Student Registration and Login:
Students can register with PRN and password and log in securely.
    \item	Practical Execution:
Students can start subject-specific practicals, compile code, get hints, and submit answers.
   
    \item	Question Management:
Admin can modify or add new questions per subject and practical.
    \item	Performance Monitoring:
The system shows progress completion bars for each subject’s practicals.
    \item	Feedback Submission:
Students can send feedback or queries to the admin using a contact form.
    \item	Session Management:
Each login is tracked, and excessive login attempts result in account lockout.

%\end{enumerate}
%\subsubsection{\textbf{User}}
%\begin{enumerate}
%    \item Give overall information to the system. 
 %    \item  View diet plan, exercise, nutritions according to BMI calculation. 
  %  \item Set their daily routine
%\end{enumerate}
\newpage
\subsection{Non Functional Requirements}
Non-functional requirements define the quality attributes of the Virtual Lab Assistant rather than the specific behaviors.
While functional requirements determine what the system should do, non-functional requirements specify how well it should perform under various conditions.
\begin{enumerate}
    \item Performance Requirement:- The system should respond within 2 seconds for major actions such as login, compile, or submit.

   \item  Availability:- The web application should be accessible 24/7, ensuring minimal downtime.
   
  \item Usability:- The interface should be user-friendly, easy to navigate, and consistent across devices.

   \item Consistency:-The interface should be user-friendly, easy to navigate, and consistent across devices.
   
  \item Reusability:- Question templates and student evaluation components can be reused for multiple subjects or academic years.

\end{enumerate} 

\section {External Interfaces}
Some External interfaces are as follows:
\begin{enumerate}
    \item User Interface
    \item Hardware Interface
    \item Software Interface

 
\subsection{User Interface}
  The proposed system has several options for users to interact with. Following are the user interfaces available:

    •	Web Application:
Provides a graphical interface for both students and administrators. Students can log in, attempt practicals, and view progress.

•	Admin Dashboard:
Allows administrators to set dates, manage subjects, and update questions.

   

   

\subsection{Hardware Interface}
  •	Compatible with desktops, laptops, and tablets.

•	Optional integration for printing student performance reports.


\subsection{Software Interface}
 •	The application supports integration with backend databases (Firebase).

•	Can connect with cloud-based storage or external APIs for future expansion.
.

\subsection{Communication Interface}
•	Uses HTTP/HTTPS for secure data transfer between the client and server.

•	Can integrate with email or SMS APIs for notifications in future updates.



\end{enumerate}






\section{Summary}
This Chapter provided a comprehensive analysis of the Virtual Lab Assistant project.
It covered requirement collection, software specifications, and external interfaces.
The study identified key functional and non-functional requirements, ensuring that the system meets academic and operational needs.
The Virtual Lab Assistant offers an innovative, scalable, and secure way to conduct online practicals efficiently while supporting both students and faculty through a seamless digital platform.
